\chapter{The story in words}\label{ch:story}

\begin{chapterplan}
What does the paper claim, and why, before any mathematics? This chapter names the players, the two returns to
information, and the chain of the main result. Later chapters prove each link.
\end{chapterplan}

\section{The question}\label{sec:story-question}

A company is for sale. One buyer, the incumbent, has already done the work that a serious offer needs. A second
buyer, the challenger, has not yet decided whether to do that work. The paper asks: can a stronger incumbent bring
the challenger into the contest?

The usual answer is no. Preparation is costly and sunk, so the challenger prepares only if its expected profit
covers the cost. A stronger incumbent lowers that profit: the challenger wins less often, and it pays more when it
wins. This force is deterrence (Fishman 1988; Hirshleifer and Png 1989).

The paper adds a force in the opposite direction. The target's stock trades while the contest is open, and the
challenger sees the price before it decides. A stronger incumbent can make the price carry news about how well the
challenger fits the target. This information effect can outweigh deterrence. Entry is the probability that the
challenger prepares.

\begin{keyidea}
A stronger incumbent has a direct effect and an indirect effect. The direct effect lowers what preparation earns.
The indirect effect makes the stock price more informative about the challenger, who reads the price before it
decides.
\end{keyidea}

\section{The set-up}\label{sec:story-setting}

In the decision interval the sale is public, the stock trades, and a further bidder can still enter.
\begin{itemize}[nosep]
\item \emph{The seller} commits to a cash second-price auction with a reserve. The winner pays the larger of the
reserve and the highest other bid. This payment is the target proceeds.
\item \emph{The incumbent} has prepared and knows its value. A stronger incumbent's value distribution moves up.
\item \emph{The challenger} has not prepared. Preparation costs zero with probability \(\rho\) and \(c_H\)
otherwise, and it shows the challenger its value. The challenger is of class \(H\) or \(L\). Before it prepares, it knows neither its class nor its value.
\item \emph{The informed investor} knows the challenger's class, but not its value. It cannot buy the target and
earns only its trading profit.
\item \emph{The liquidity trader} sells one unit, buys one unit, or does nothing. \emph{The market makers} see the
order only, and set the price equal to the expected payoff of a share.
\end{itemize}
The class ranking says two things: at every level, class \(H\) is at least as likely as class \(L\) to value the
target above that level, and every value meets the reserve. The prior is the probability of class \(H\) before any
trade. Trade takes place in one order slot. The informed investor fills the slot with some
probability, and a liquidity trader fills it otherwise. A purchase or a sale costs the investor a trading cost.
Nobody sees who filled the slot. The entry floor \(\rho\) is the least possible
entry, because a zero-cost challenger prepares after every price.

In order: the rule; the class to the investor; the order; the price; the challenger's cost and decision; the bids.

\begin{watchout}
The challenger sees the price and its cost, but not the order. The investor sees the class, but not the value. The
price carries news only through what the market makers infer from the order.
\end{watchout}

\section{The two returns to information}\label{sec:story-returns}

The challenger's gross acquisition profit is what it expects to keep from a win, before its cost. If class \(H\)
expects more than class \(L\), gross profit rises with the challenger's belief. The preparation threshold is the
belief at which gross profit equals \(c_H\). The target-payoff spread \(\DT\) is the gap between the classes in the
target's expected proceeds. It is what the investor's knowledge of the class is worth.

Proposition~1 says that a stronger incumbent (a) never raises the challenger's gross profit at any belief, and
(b) under the class ranking, never narrows the spread. One picture gives both. Cut the gap between the classes into
thin slices, one for each level of value above the reserve. In a second-price auction, the slices below the
incumbent's value go to the target, and the slices above it stay with the challenger. A stronger incumbent cuts
further to the right on average. So the challenger keeps less, and the spread rises.
Both statements are weak: a strict change needs a condition on the slices (Chapter~\ref{ch:auction}).

\section{The main result as a chain}\label{sec:story-chain}

Proposition~2 compares three economies: a weak incumbent, a strong one, and a very strong one. All else is the same. Table~\ref{tab:story-three} gives the outcomes.

\begin{table}[ht]
\centering\footnotesize
\caption{The three economies of Proposition~2, in words.}\label{tab:story-three}
\begin{tabular}{@{}>{\raggedright\arraybackslash}p{2.3cm}>{\raggedright\arraybackslash}p{3.6cm}>{\raggedright\arraybackslash}p{4.4cm}>{\raggedright\arraybackslash}p{3.6cm}@{}}
\toprule
 & Weak incumbent & Strong incumbent & Very strong incumbent\\
\midrule
The spread & Below the trading cost & Wide: a correct order covers the trading cost & At least as wide\\
The investor & Does not trade & Class \(H\) buys, class \(L\) sells & The same orders\\
The price & Silent: belief at the prior & Informative & Informative\\
Gross profit & At the prior, below \(c_H\) & At the top belief, above \(c_H\) & Top belief, below \(c_H\)\\
Entry & The floor \(\rho\) & Above the floor & Back at the floor \(\rho\)\\
\bottomrule
\end{tabular}
\end{table}

The chain has four links.
\begin{enumerate}[nosep]
\item A wider spread makes informed trade pay. A correct order earns at most the spread. Against the weak incumbent
this is below the trading cost. Against the strong incumbent a correct order earns at least a fixed share of the
spread, which exceeds the trading cost.
\item Informed trade makes the price informative. A purchase is likelier in class \(H\), so it raises the belief.
The challenger reads the belief from the price.
\item After a purchase the belief equals the top posterior. Against the strong incumbent, gross profit there is above
\(c_H\), so the purchase recruits a high-cost challenger.
\item Entry is one after a purchase and \(\rho\) otherwise, so entry rises above the floor.
\end{enumerate}
High-class ownership, the probability that a class-\(H\) challenger acquires the target, also rises when class \(H\) outbids the strong incumbent for sure.

Against a very strong incumbent the spread is at least as wide, so the investor still trades. But gross profit has
fallen at every belief. Even the top posterior leaves it below \(c_H\), so no price recruits, and entry returns to
the floor.

\begin{watchout}
The three economies are separate, and the equilibrium is solved again in each one. The result compares targets that
differ only in the incumbent. It does not follow one target as its incumbent grows.
\end{watchout}

The hypotheses are three windows: a high-cost window (H1), a trading-cost window (H2) and an exclusion condition
(H3). They hold together on a nonempty open set (Chapters~\ref{ch:ccc} and~\ref{ch:example}).

\section{Why the price never reveals everything}\label{sec:story-noise}

A liquidity trader can place any order in either class, so every order has positive probability in both classes.
So every posterior, and every belief at a price, lies between a bottom
and a top posterior. For the order of no trade this needs the chance of a liquidity purchase to be at most one
third.

The investor earns the entry that its order brings, times the spread, times the remaining doubt about the class. Entry is at least \(\rho\) and the doubt has a floor, so a correct order earns at least a fixed share of the spread: the trading bound. The price reveals the market makers' belief, but never the class.

\section{Two controls}\label{sec:story-controls}

The first control holds the information in the price fixed. A stronger incumbent then lowers gross profit at every
belief, so entry can only fall. In the example, with the investor's orders held at their informative level, entry
stays flat from the weak to the strong incumbent. The second control hides
the price. The challenger then decides on its prior, which lies below the threshold in the example, so entry stays
at the floor although the investor trades.

So entry rises only when trading creates information and the challenger reads it from the price.

\section{Quiet markets, research, data and the seller}\label{sec:story-markets}

When the spread is positive and class \(H\) earns more, the market depends on the auction only through the spread and the threshold. Proposition~3 shows that between the weak and the strong incumbent a quiet market and an informative market can coexist against the same incumbent. Each state
confirms itself. When the investor pays for research (Proposition~4), research and entry rise with the spread until research reaches its cap, in the equilibrium where a purchase after a success recruits. The rise in entry survives
other cost laws, more classes and private signals under stated conditions (Chapter~\ref{ch:robustness}).

Prices and order counts can test the story (Chapter~\ref{ch:empirical}). For the seller (Chapter~\ref{ch:welfare}), under the hypotheses of Proposition~2, access to prices raises
expected target proceeds and net acquisition surplus against the strong incumbent. 
\begin{chaptersummary}
A stronger incumbent never raises the challenger's gross profit (deterrence) and, under the class ranking, never narrows the spread. A wider spread makes informed trade pay, trade makes the price
informative, and a purchase can lift the belief past the preparation threshold, so entry rises. Against a very
strong incumbent even the top belief leaves profit below the high cost, and entry returns to the floor.
\end{chaptersummary}
